\documentclass{article}
\usepackage{../../shared/omnicourse}
\input{../../shared/meta.tex}

\renewcommand{\docshortname}{Software Development Plan}
\hypersetup{pdftitle={OmniCourse -- Software Development Plan v2.0}}

\begin{document}

\makecover{Software Development Plan}{2.0}{1 November 2026}{ESE-PA2-SDP-2.0}{Project Assignment 2 (PA2)}

\pagenumbering{roman}
\thispagestyle{empty}

\section*{Revision history}
\vspace{-0.3em}
{\small
\begin{xltabular}{\dimexpr\textwidth-2\tabcolsep\relax}{@{}p{2.6cm}p{1.8cm}Yp{4.2cm}@{}}
\hcell{Date} & \hcell{Version} & \hcell{Description} & \hcell{Author}\\\hdrline
\endhead
18/Oct/2026 & 1.0 & Initial version submitted for PA1 & Lê Hữu Hoà\\
1/Nov/2026 & 2.0 & Second submission for PA2: updated schedule, expanded resource allocation and monitoring sections. Revisions highlighted in blue. & Phan Bình Nguyên Lâm, Lê Hữu Hoà\\
\end{xltabular}}
\par\vspace{1.0cm}

{\hypersetup{linkcolor=ocNavy}\tableofcontents}

\clearpage
\pagenumbering{arabic}

\section{Introduction}
This Software Development Plan describes how the OmniCourse team of five students organises, schedules and controls the
development of the product in the course project. The project runs for twelve weeks, from 5 October to 13 December 2026,
organised as five two-week sprints inside three RUP phases. The document defines the purpose and scope, the team
organisation and roles, the phase and iteration plan, the planned releases, the schedule, the resourcing, and the
reporting and configuration management processes. It is a first draft written in the Inception phase. The project
estimates and the risk register are left empty and are completed in the next assignment, as PA1 requires.

\section{Project Overview}
\subsection{Project Purpose, Scope, and Objectives}
OmniCourse is a public web application that gives a HCMUS student one page for everything school related: degree audit,
curriculum explorer, semester roadmap planner, an AI planner and advisor, and the exact links and steps for the school
tasks the application does not perform itself. The first release covers the Advanced Program in Computer Science (APCS);
the model is program-agnostic, so other HCMUS majors and later other Vietnamese universities are added as configuration.

The objectives of the project are:

\begin{enumerate}
  \item a deployed, working web application that covers the APCS curriculum;
  \item a deterministic rule engine that enforces prerequisites, offering terms, requirement coverage and the credit caps
        of the program, so that no academic rule is decided by a language model;
  \item an AI plan generator that turns a plain-language goal into a roadmap validated by that engine and explains every
        placement;
  \item a curriculum data model that can be reused for another program or another school without changing the engine.
\end{enumerate}

Out of scope for this course: native mobile applications, automation of the school registration portal, storage of portal
credentials, and the optional modules (tuition, conduct score, paperwork, community).

\subsection{Assumptions and Constraints}
The plan rests on the following assumptions and constraints:

\begin{itemize}
  \item Fixed schedule: five sprints of two weeks, 5/10/2026 to 13/12/2026, aligned with the course plan agreed in PA0.
        No extension is possible.
  \item Team of five students, no member added or removed during the project.
  \item Zero budget. Only free tiers are used: hosting, CI minutes and the model API.
  \item Expected effort is at least 10 hours per member per week, about 50 hours per week for the team.
  \item All five members are full-time students, so lecture and examination weeks reduce capacity.
  \item The process is fixed by the course: three RUP phases, five Scrum sprints, and the tools approved in PA0
        (Moodle, Discord, GitHub).
  \item No official API exists for the HCMUS curriculum or the student record, so curriculum data is entered and verified
        by the team, and the transcript is a file the student uploads.
  \item The academic rules follow Thông tư 08/2021/TT-BGDĐT and the published HCMUS credit caps (14 to 25 credits in a
        main term, 12 in summer). Personal data follows Luật 91/2025/QH15.
  \item The application never stores a portal password and never automates a registration on the student's behalf.
\end{itemize}

\subsection{Project Deliverables}
The artifacts below are produced during the project. Documents are stored in the GitHub repository under \texttt{/docs}
and submitted in \texttt{/report}, as PA0 requires.

{\small
\begin{xltabular}{\dimexpr\textwidth-2\tabcolsep\relax}{@{}Yp{2.5cm}p{2.2cm}@{}}
\hcell{Artifact} & \hcell{Target date} & \hcell{Assignment}\\\hdrline
\endhead
Project plan v1.0 & 18/10/2026 & PA1\\
Vision document v1.0 & 18/10/2026 & PA1\\
UI sketches of the main screens (3 to 6) & 18/10/2026 & PA1\\
Weekly reports, Sprint 1 (2 reports) & 18/10/2026 & PA1\\
Project plan v2.0 with TA feedback applied & 1/11/2026 & PA2\\
Use-case model and use-case specifications & 1/11/2026 & PA2\\
Vision document v2.0 & 1/11/2026 & PA2\\
Software architecture document, initial & 1/11/2026 & PA2\\
Detailed design: class and sequence models & 15/11/2026 & PA3\\
Test plan and test cases & 15/11/2026 & PA3\\
Build B1 deployed on staging & 15/11/2026 & PA3\\
Build B2, working software & 29/11/2026 & PA4\\
Build B3 release candidate and test report & 13/12/2026 & PA5\\
\end{xltabular}}

\section{Project Organization}
\subsection{Organizational Structure}
The team has a flat structure of five members. The project lead coordinates the plan, the schedule and the sharing of work,
and is the single contact with the instructor and the TAs, who review and grade each assignment and may request changes to
the scope. Two technical leads own the frontend and the backend workstreams and review the pull requests of their side; the
two remaining members carry the testing duties. The structure is deliberately small: the stack, the data model and the
scope of each sprint are decided together in the weekly meeting, and the project lead breaks a tie.

\begin{itemize}
  \item Project Lead, the contact with the instructor and the TAs.
  \item Frontend Lead and Backend Lead, one workstream each and the review of its pull requests.
  \item Two members responsible for testing in each sprint, in addition to their development work.
\end{itemize}

\subsection{Roles and Responsibilities}
All five members write code, so every member is a Developer. The three leads carry coordination duties on top of that, and
the two remaining members own the testing of each sprint. Roles are assigned per sprint, not fixed for the whole project,
so each member rotates through analysis, implementation and testing across the five sprints. The table gives the standing
assignment and the duty that belongs to the role.

{\small
\begin{xltabular}{\dimexpr\textwidth-2\tabcolsep\relax}{@{}p{3.9cm}p{3.7cm}Y@{}}
\hcell{Member} & \hcell{Role} & \hcell{Main responsibilities in this project}\\\hdrline
\endhead
\textbf{Phan Bình Nguyên Lâm} & Project Lead, Developer &
Owns the plan and the schedule, chairs the weekly meeting, keeps the workload balanced, tracks progress and reports to the
instructor and the TAs. Implements backend features.\\
\textbf{Lê Hữu Hoà} & Frontend Lead, Developer &
Owns the Vue application: components, routing, state and consistency with the UI sketches. Reviews frontend pull requests.
Implements frontend features.\\
\textbf{Võ Nguyễn Minh Triết} & Backend Lead, Developer &
Owns the NestJS API, the PostgreSQL schema and the rule engine. Reviews backend pull requests. Implements backend
features.\\
\textbf{Nguyễn Anh Dũng} & Developer, Tester &
Implements features. Owns the test plan and the test cases of the sprint and runs the regression suite on every build.\\
\textbf{Nguyễn Hữu Quốc Thắng} & Developer, Tester &
Implements features. Maintains the automated test suite (Vitest) and checks each deployed build against the acceptance
criteria of the sprint.\\
\end{xltabular}}

\section{Management Process}
\subsection{Project Estimates}
\begin{ocplain}{Not part of this version}
The estimates of cost and schedule are not required in PA1. The section is completed in PA2, once the requirements
baseline is approved and the size of the remaining work is known.
\end{ocplain}

\subsection{Project Plan}
This section contains the schedule and the resources of the project: the phases and iterations, the planned releases, the
schedule of each sprint with the tasks and their owners, and the resourcing.

\subsubsection{Phase and Iteration Plan}
The project has three RUP phases and five two-week iterations. The boundaries follow the overall plan agreed in PA0, so
the plan, the vision document and the schedule of the assignment stay consistent.

{\small
\begin{xltabular}{\dimexpr\textwidth-2\tabcolsep\relax}{@{}p{2.2cm}p{1.4cm}p{1.7cm}p{1.7cm}Y@{}}
\hcell{Phase} & \hcell{Iteration} & \hcell{Start} & \hcell{End} & \hcell{Overall objectives}\\\hdrline
\endhead
Inception & Sprint 1 & 5/10/2026 & 18/10/2026 &
Project proposal, group and tools set up, first versions of the project plan and the vision document, UI sketches.\\
Elaboration & Sprint 2 & 19/10/2026 & 1/11/2026 &
Requirements baseline, use cases, initial architecture and database schema, running skeleton (build B0).\\
Elaboration & Sprint 3 & 2/11/2026 & 15/11/2026 &
Detailed design, test plan and test cases, curriculum data model, degree audit and curriculum explorer (build B1).\\
Construction & Sprint 4 & 16/11/2026 & 29/11/2026 &
Roadmap planner with rule validation, transcript import, notifications (build B2).\\
Construction & Sprint 5 & 30/11/2026 & 13/12/2026 &
AI plan generator and advisor, quicklinks, performance and security pass, regression, release candidate (build B3).\\
\end{xltabular}}

Milestones: M1 documents of PA1 approved, 18/10; M2 requirements baseline, 1/11; M3 architecture, design and test plan
approved, 15/11; M4 core planning features working, 29/11; M5 release candidate and final report, 13/12.

\subsubsection{Releases}
Four builds are produced for testing. Each build is deployed to the staging environment from the main branch through the
continuous integration pipeline. On every build the testers run the regression suite plus the acceptance criteria of the
sprint. Only build B3 is a release.

{\small
\begin{xltabular}{\dimexpr\textwidth-2\tabcolsep\relax}{@{}p{1.1cm}p{4.2cm}p{3.2cm}Y@{}}
\hcell{Build} & \hcell{Date} & \hcell{Type} & \hcell{Content}\\\hdrline
\endhead
B0 & 1/11/2026, end of Sprint 2 & Demo, walking skeleton &
Deployment pipeline, database schema, sign-in, one working screen.\\
B1 & 15/11/2026, end of Sprint 3 & Demo, beta &
Degree audit dashboard and curriculum explorer on the APCS curriculum.\\
B2 & 29/11/2026, end of Sprint 4 & Beta &
Semester roadmap planner with live rule validation, transcript import, notifications.\\
B3 & 13/12/2026, end of Sprint 5 & Release candidate &
AI plan generator and advisor, quicklinks and guides, full regression.\\
\end{xltabular}}

\subsubsection{Project Schedule}
The schedule below is the one the team works to. Tasks are kept in the project Discord server, one message per task, moved
by its owner through the states Pending, Doing and Done. This table is copied from that list at the end of each sprint, so
progress can be checked against the dates. The schedule may change after the review of PA1.

{\small
\begin{xltabular}{\dimexpr\textwidth-2\tabcolsep\relax}{@{}p{1.4cm}p{2.0cm}Yp{3.0cm}@{}}
\hcell{Sprint} & \hcell{Dates} & \hcell{Key tasks} & \hcell{Owner}\\\hdrline
\endhead
Sprint 1 & 5/10 -- 18/10 &
Project proposal and 5 user interviews; group registration; tools setup; stack decision; project plan; vision document;
UI sketches; 2 weekly reports. & All, coordinated by Lâm\\
Sprint 2 & 19/10 -- 1/11 &
Use-case model and specifications; requirements baseline; initial architecture and database schema; continuous
integration; build B0 deployed. & Lâm, Triết, Thắng\\
Sprint 3 & 2/11 -- 15/11 &
Detailed design; test plan and test cases; curriculum data model for APCS; degree audit; curriculum explorer; build B1.
& Hòa, Triết, Dũng\\
Sprint 4 & 16/11 -- 29/11 &
Roadmap planner with validation; transcript import; notifications; build B2; regression. & All\\
Sprint 5 & 30/11 -- 13/12 &
AI plan generator and advisor; quicklinks and guides; performance and security pass; full regression; release candidate
B3; final report. & All\\
\end{xltabular}}

\subsubsection{Project Resourcing}
Five members, each committing at least 10 hours per week, about 50 hours per week and roughly 600 hours over the twelve
weeks. The work is spread so that no sprint depends on a single person: every task has an owner and a reviewer.

\begin{itemize}
  \item Source control and review: GitHub, one repository, protected main branch.
  \item Communication and task tracking: one Discord server, one channel per concern, one message per task.
  \item Development: NestJS and Vue run locally; PostgreSQL 18 locally with seeded curriculum data.
  \item Continuous integration and hosting: GitHub Actions and a free-tier cloud host for staging.
  \item Documents: Markdown in \texttt{/docs}, exported to \texttt{.pdf} for submission.
  \item No paid licence is needed, which keeps the project inside the zero-budget constraint.
\end{itemize}

\subsection{Project Monitoring and Control}
\subsubsection{Reporting}
Status is reported in three ways. First, a weekly report every Sunday, written on the course template and stored in
\texttt{/docs/management}: it lists what was completed, what slipped, the issues and the goals of the following week.
Second, a weekly meeting of about thirty minutes at the start of each week, where the report is closed, the priorities of
the sprint are confirmed and any blocker is assigned to a person. Third, continuous status in the project Discord server:
every task is a message in the channel and its owner moves it through the states Pending, Doing and Done on the day it
changes, so the state of the sprint can be read without asking anyone. Two weekly reports are produced during Sprint 1
and submitted with PA1.

\begin{itemize}
  \item Weekly report to the TAs: what was planned, what was done, the percentage of each task, the issues and the goals of
        the next week.
  \item Weekly meeting: thirty minutes, decisions recorded in the channel right after the meeting.
  \item Daily updates in Discord: task state, blockers, and the pull request number once the work is opened for review.
\end{itemize}

\subsubsection{Risk Management}
\begin{ocplain}{Not part of this version}
The risk register is left empty in PA1, as the assignment instructs. The table of risks, probabilities and mitigation
strategies is completed in PA2, where the technical risks of the rule engine and of the curriculum data become measurable.
\end{ocplain}

\subsubsection{Configuration Management}
Source code and documents live in one GitHub repository owned by the team, with the folder layout required by PA0:
\texttt{/src} for the code, \texttt{/docs/management} for plans and reports, \texttt{/docs/requirements} for the vision
document and the use cases, \texttt{/docs/analysis and design} for the architecture, the models and the UI design,
\texttt{/docs/test} for the test plan, the test cases and the test reports, and \texttt{/report} for the submissions.
Documents are written in Markdown, versioned with the code, and exported to \texttt{.pdf} only for submission.

\begin{itemize}
  \item No one pushes to main. Every change arrives as a pull request.
  \item A pull request merges only when the continuous integration pipeline passes, every commit is signed, and at least
        one other member has reviewed and approved it.
  \item After approval the author merges the pull request and performs the deployment.
  \item Review duty follows the workstream: the frontend lead reviews frontend work, the backend lead reviews backend work,
        and the project lead reviews changes to the data model and the rules.
  \item Builds B0 to B3 are tagged in the repository, so the version that was tested is always identifiable.
\end{itemize}

The technology stack of the project, all at the current stable release:

{\small
\begin{xltabular}{\dimexpr\textwidth-2\tabcolsep\relax}{@{}p{3.0cm}Yp{4.0cm}@{}}
\hcell{Layer} & \hcell{Technology} & \hcell{Version}\\\hdrline
\endhead
Database & PostgreSQL & 18.6\\
Backend & NestJS, with the Nest CLI and Swagger & 12.1.3 / 12.0.8 / 12.0.2\\
Frontend & Vue & 3.5.44\\
API & REST, described by OpenAPI & 3.2.0\\
Tests & Vitest & 5.0.4\\
Build & Vite & 8.3.4\\
\end{xltabular}}

\end{document}
